
We proceed with the three steps of the modeling process.
\begin{description}
\item[Decision variables] The decision variables are the amounts in euros invested
  for
\begin{itemize}
\item   state bonds: $x_{sb}$,
\item real estate loans: $x_{re}$,
\item car loans:  $x_{c\ell}$, and,
 \item scholarship loans: $x_{s\ell}$.
\end{itemize}
\item[Objective function] As the company wants to maximize its profit,
  the objective function must provide the profit as a function of the
  decision variables:
  \[
f:\R^4 \to \R: f(x_{sb}, x_{re}, x_{c\ell},x_{s\ell}).
  \]
The profit of the investment fund for the following year is calculated
based on  the interest rates using the following formula:
\[
f(x_{sb}, x_{re}, x_{c\ell},x_{s\ell}) = 0.04 x_{sb} + 0.06 x_{re} + 0.08 x_{cl} + 0.09 x_{s\ell}. 
\]
Therefore, we can write the problem as 
\[
	\max_{x_{sb}, x_{re}, x_{c\ell},x_{s\ell}} 0.04 x_{sb} + 0.06 x_{re} + 0.08 x_{cl} + 0.09 x_{s\ell}. 
\]
\item[Constraints]
The total amount to invest is 2.5 M\euro.  This is modeled using the following constraint:
\[ 	x_{sb} + x_{re} + x_{c\ell} + x_{s\ell}  = 2'500'000 \]

Each restriction is modeled as follows:
\begin{itemize}
	\item the total amount invested in car and scholarship  loans
          ($	x_{c\ell} + x_{s\ell}$) must not exceed twice the
          amount invested in bonds ($2x_{sb}$):
	\[ 	x_{c\ell} + x_{s\ell}  \leq 2x_{sb}. \]
	\item the amount invested in car loans ($x_{c\ell}$) must be
          larger or equal than the
          amount invested in  scholarship loans ($x_{s\ell}$):
	\[ 	x_{s\ell}  \leq x_{c\ell}. \]
	\item the investment in car loans ($x_{c\ell}$) should not exceed the
          investment in real estate loans ($x_{re}$) by more than $70\%$:
	\[ 	x_{c\ell}  \leq 1.7x_{re}. \]
\end{itemize}  
Finally, we must impose all the decision variables to be non negative:
\[
	x_{sb}, x_{re}, x_{c\ell}, x_{s\ell} \geq 0.
\]
\end{description}



Putting everything together, we obtain the following optimization problem:
\[
\max 0.04 x_{sb} + 0.06 x_{re} + 0.08 x_{c\ell} + 0.09 x_{s\ell} 
\]
subject to
\begin{align*}
	x_{sb} + x_{re} + x_{c\ell} + x_{s\ell}  &= 2'500'000, \\
	x_{c\ell} + x_{s\ell}  &\leq 2x_{sb}, \\
	x_{s\ell}  &\leq x_{c\ell}, \\
	x_{c\ell} &\leq 1.7x_{re}, \\
	x_{sb}, x_{re}, x_{c\ell}, x_{s\ell} &\geq 0.
\end{align*}

For your information, the optimal solution (rounded to 1/10 of euros) is:
\begin{itemize}
\item State bonds: 696'721.3 euros
\item Real estate loans: 409'836.1 euros
\item Car loans: 696'721.3 euros
\item Scholarship loans: 696'721.3 euros
\end{itemize}
for a total profit of 170'902.6 euros

The constraints are verified:
\begin{itemize}
\item Constraint 1: $696'721.3 + 409'836.1 + 696'721.3 + 696'721.3 = 2'500'000.0$,
\item Constraint 2: $696'721.3 + 696'721.3 = 1'393'442.6  2 * 696'721.3 = 1'393'442.6$,
\item Constraint 3: $696'721.3 <= 696'721.3$,
\item Constraint 4: $696'721.3 <= 1.7 \cdot 409'836.1 = 696'721.3$.
\end{itemize}


